\section*{Bab \ref{ch:b2:affine} --- Ruang Afin}

\begin{solution}{exo:b2:affine:1}
$\{x + y + z = 1\}$: bidang \omterm{def:b2:affine:subspace}{afin}, berarah bidang vektor $\{x
+ y + z = 0\}$, dan berdimensi $2$. Menambahkan $x - z = 3$: garis \omterm{def:b2:affine:subspace}{afin}
(sebab dua persamaannya bebas), berarah $\{x + y + z = 0,\ x = z\}
= \operatorname{Vect}\bigl((1, -2, 1)\bigr)$ dan berdimensi $1$. $\{x^2
+ y^2 = 1\}$: sebuah silinder --- yang tak stabil terhadap \omterm{def:b2:affine:barycenter}{barisentrum} (sebab
titik tengah $(1,0,0)$ dan $(-1,0,0)$ adalah titik asalnya, di luar
silindernya): jadi tak \omterm{def:b2:affine:subspace}{afin}. Adapun sistem serasi $MX = B$: yakni subruang
\omterm{def:b2:affine:subspace}{afin} $X_0 + \ker M$ berdimensi $\dim\ker M$, seperti diingatkan pada
\cref{def:b2:affine:subspace}.
\end{solution}

\begin{solution}{exo:b2:affine:2}
$C = \operatorname{bar}(A, 1-t; B, t)$ berarti $\vect{AC} =
t\,\vect{AB}$: jadi keberadaan $t$ semacam itu persis merupakan kebergantungan
$\vect{AC}$ dengan $\vect{AB} \neq 0$, yakni kesegarisan. Adapun letaknya: $t
= \frac12$: titik tengahnya; $t = 2$: melewati $B$, sejauh jarak $B$ dari
titik itu (sebab $\vect{AC} = 2\vect{AB}$); dan $t = -1$: cerminan $B$
terhadap $A$.
\end{solution}

\begin{solution}{exo:b2:affine:3}
Matriks $M$ adalah matriks proyeksi ke $\operatorname{Vect}(1,1)$ sepanjang
$(1,-1)$ (periksa $M^2 = M$). Peta $f$: yakni $\{MX + C\} = C +
\operatorname{im} M$, jadi garis \omterm{def:b2:affine:subspace}{afin} lewat $(1,0)$ yang diarahkan
$(1,1)$. Titik tetapnya: $X = MX + C$, yakni $(I - M)X = C$; tetapi $C =
(1, 0)^{\mathsf T}$ dan $\operatorname{im}(I - M) =
\operatorname{Vect}(1,-1)$; apakah $(1,0)$ ada di dalamnya? Sebab $(1, 0) =
\alpha(1,-1)$ memaksa $\alpha = 1$ dan $0 = -1$: jadi tidak. Sehingga tak ada titik
tetap. Dan
\[
f(f(X)) = M(MX + C) + C = MX + MC + C = f(X) + MC,
\qquad MC = \tfrac12(1,1)^{\mathsf T} :
\]
jadi $f\circ f$ adalah $f$ yang disusul penggeseran sepanjang garis petanya
--- sehingga $f$ merupakan ``proyeksi geser'': yakni proyeksi ke garisnya
yang dikomposisikan dengan sebuah geseran.
\end{solution}

\begin{solution}{exo:b2:affine:4}
$I = \operatorname{bar}(B, 2; C, 1)$ (sebab $\vect{BI} =
\frac13\vect{BC}$ menempatkan $I$ lebih dekat ke $B$: jadi bobot $2$ pada $B$
dan $1$ pada $C$ --- periksa: $\vect{BI} = \frac{1}{3}\vect{BC}$).
Demikian pula $J = \operatorname{bar}(C, 2; A, 1)$ dan $K =
\operatorname{bar}(A, 2; B, 1)$. Menjumlahkan ketiga sistem terbobotnya,
\omterm{def:b2:affine:barycenter}{barisentrum} $(I, 1; J, 1; K, 1)$ (yang masing-masing berbobot total
$3$, jadi gantilah $I$ dengan sistemnya, dan seterusnya) adalah
\[
\operatorname{bar}\bigl(A, 1 + 2;\ B, 2 + 1;\ C, 1 + 2\bigr)
= \operatorname{bar}(A, 1; B, 1; C, 1) = G ,
\]
yakni titik berat $ABC$: jadi segitiga $IJK$ punya titik berat yang sama ---
dan itu dapat diramalkan, karena konstruksinya memperlakukan $A, B, C$ secara \omterm{def:b2:structures:generated}{siklik}
sedangkan titik beratnya adalah satu-satunya titik tetap kesimetrian \omterm{def:b2:structures:generated}{siklik}
bobotnya.
\end{solution}

\begin{solution}{exo:b2:affine:5}
Tetapkan $g(u) = f(O + u) - f(O)$ (sambil bekerja di $\R^n$ yang divektorkan di
$O$), dengan $g(0) = 0$.

\emph{Keaditifannya:} sebab $\frac{(O + u) + (O + v)}{2} = O + \frac{u +
v}{2}$, jadi pemeliharaan titik tengahnya memberi $g\bigl(\frac{u+v}{2}\bigr)
= \frac{g(u) + g(v)}{2}$; lalu dengan $v = 0$: $g(u/2) = g(u)/2$;
sehingga bila digabung, $g(u + v) = 2g\bigl(\frac{u+v}{2}\bigr) = g(u) + g(v)$.

\emph{Kehomogenan atas $\Q$:} keaditifannya memberi $g(nu) = ng(u)$ (untuk $n \in
\N$, secara induktif), lalu $g(-u) = -g(u)$ (dengan menjumlahkan), lalu $g(\frac pq u) =
\frac pq g(u)$ (terapkan $q$, lalu pakailah keinjektifan penskalaannya).

\emph{Kehomogenan atas $\R$:} untuk $t \in \R$, ambillah bilangan rasional $t_n \to
t$: maka $g(t_nu) = t_ng(u)$, lalu kekontinuan $g$ (yang diwarisi dari $f$)
melewatkan limitnya: sehingga $g(tu) = tg(u)$. Karena itu $g$ bersifat linear dan $f =
f(O) + g$: jadi \omterm{def:b2:affine:subspace}{afin}.
\end{solution}

\begin{solution}{exo:b2:affine:6}
Bagian linearnya $\varphi(x,y) = (y, x)$: yakni pencerminan pada diagonal
$y = x$ (yang ortogonal, berdeterminan $-1$). Titik tetapnya: $(x, y) = (y
+ 1, x - 1)$ menyusut menjadi satu persamaan tunggal $y = x - 1$ (sebab kedua
komponennya setara): jadi setiap titik garis $y = x - 1$ bersifat
tetap. Sehingga $f$ menetapkan garis itu \omterm{def:b2:funcseq:def}{titik demi titik}: jadi $f$
adalah \emph{pencerminan} pada sumbu itu (yakni isometri dengan sebuah garis titik
tetap dan bagian linear berupa pencerminan). Secara konsisten, $f \circ
f(x,y) = f(y+1, x-1) = (x - 1 + 1, y + 1 - 1) = (x, y)$: jadi sebuah
involusi, sebagaimana mestinya bagi sebuah pencerminan.
\end{solution}

\begin{solution}{exo:b2:affine:7}
Adapun $n + 1$ vektor $\vect{A_1A_i}$ (untuk $i \geq 2$) bergantungan pada
dimensi $n$: jadi ada $\mu_i$ yang tak semuanya nol dengan $\sum_{i\geq2}
\mu_i\vect{A_1A_i} = 0$; lalu tetapkan $\mu_1 = -\sum_{i \geq 2}\mu_i$, sehingga
$\sum_{i}\mu_i = 0$ dan $\sum_i \mu_i\,\vect{OA_i} = 0$ untuk setiap
$O$, dengan tak semua $\mu_i$ nol. Pecahlah indeksnya: $P = \{i : \mu_i >
0\}$ dan $N = \{i : \mu_i < 0\}$, yang keduanya tak kosong (sebab $\mu_i$ berjumlah
nol dan tak semuanya nol). Dengan $s = \sum_{i\in P}\mu_i =
-\sum_{i \in N}\mu_i > 0$:
\[
\operatorname{bar}\bigl(A_i, \tfrac{\mu_i}{s}\bigr)_{i \in P}
= \operatorname{bar}\bigl(A_i, \tfrac{-\mu_i}{s}\bigr)_{i \in N},
\]
(sebab kedua ruasnya sama dengan titik $X$ yang $\vect{OX} = \frac1s\sum_{i\in
P}\mu_i\vect{OA_i}$, menurut hubungannya): jadi sebuah titik bersama kedua
\omterm{def:b2:affine:convex}{bungkus cembungnya}, dengan kelompok indeks yang saling lepas.
\end{solution}

\begin{solution}{exo:b2:affine:8}
\emph{Peta = titik tetapnya:} sebab untuk $Y = f(X)$ berlaku $f(Y) = f(f(X)) =
f(X) = Y$: jadi setiap titik petanya bersifat tetap; dan sebaliknya titik tetapnya adalah
peta. Karena itu $\mathcal{F} = \operatorname{im} f =
\operatorname{Fix}(f)$ tak kosong, dan ia sebuah subruang \omterm{def:b2:affine:subspace}{afin}
(sebab peta sebuah \omterm{def:b2:affine:subspace}{pemetaan afin}), dengan arah $\operatorname{im}\vec f$.

\emph{Struktur proyeksinya:} $\vec f$ bersifat idempoten (sebab $\vec{f\circ
f} = \vec f^{\,2} = \vec f$), jadi $E = \operatorname{im}\vec f
\oplus \ker \vec f$ (\cref{ex:b2:reduction:projections}). Untuk sembarang
titik $X$, tinjaulah vektor $\vect{f(X)\,X}$; lalu dengan menerapkan $\vec f$:
\[
\vec f\bigl(\vect{f(X)\,X}\bigr) = \vect{f(f(X))\,f(X)} = 0
\qquad (f \circ f = f),
\]
jadi $\vect{f(X)\,X} \in \ker\vec f$. Karena itu $X = f(X) +
\vect{f(X)X}$ memamerkan $X$ sebagai titik $\mathcal{F}$ yang digeser
sebuah vektor $\ker\vec f$: sehingga $f$ persis merupakan proyeksi ke
$\mathcal{F}$ sepanjang $\ker\vec f$. Sebaliknya proyeksi semacam itu
jelas memenuhi $f \circ f = f$.
\end{solution}

\begin{solution}{exo:b2:affine:9}
Misalkan $M = \operatorname{bar}(B, 2;\ C, 3)$, yang berbobot total $5$.
Keasosiatifannya memberi $G = \operatorname{bar}(A, 1;\ M, 5)$, jadi
$\vect{AG} = \frac56\,\vect{AM}$: sehingga $G$ terletak pada ruas
$\intcc AM$ pada lima perenamnya dari $A$. Karena $A \notin
(BC)$, garis $(AG) = (AM)$ memotong $(BC)$ di titik tunggal
$M$, dengan $\vect{BM} = \frac35\,\vect{BC}$. Demikian pula, dengan $N =
\operatorname{bar}(C, 3;\ A, 1)$ (yang berbobot total $4$ dan $\vect{CN}
= \frac14\,\vect{CA}$), keasosiatifannya memberi $G =
\operatorname{bar}(B, 2;\ N, 4)$: jadi garis $(BG)$ memotong $(CA)$
di $N$, dan $\vect{BG} = \frac46\,\vect{BN} =
\frac23\,\vect{BN}$.
\end{solution}

\begin{solution}{exo:b2:affine:10}
Vektorkanlah di sebuah titik asal $O$ lalu tulislah titiknya sebagai vektor:
$h_{\Omega, \lambda}(x) = \omega + \lambda(x - \omega)$ dengan
$\omega = \vect{O\Omega}$. Komposisi $g = h_{\Omega',
\mu} \circ h_{\Omega, \lambda}$ bersifat \omterm{def:b2:affine:subspace}{afin} dengan bagian linear
$\mu\lambda\,\mathrm{id}$. Bila $\lambda\mu \neq 1$: maka $1 \notin
\operatorname{Sp}(\lambda\mu\,\mathrm{id})$, jadi
\cref{prop:b2:affine:fixedpoint} menghasilkan titik tetap tunggal
$\Omega''$ dan, setelah divektorkan di sana, $g =
\lambda\mu\,\mathrm{id}$: yakni homoteti $h_{\Omega'',
\lambda\mu}$. Sedangkan bila $\lambda\mu = 1$ bagian linearnya adalah
identitasnya, jadi $g$ merupakan penggeseran; dan setelah diuraikan,
\[
g(x) = \omega' + \mu\bigl(\omega + \lambda(x - \omega) -
\omega'\bigr) = x + (1 - \mu)\,\omega' + \mu(1 -
\lambda)\,\omega .
\]
Untuk $\lambda = \mu = -1$ (yakni pencerminan titik) vektornya adalah
$2\omega' - 2\omega = 2\,\vect{\Omega\Omega'}$: jadi
komposisi pencerminan titik pada $\Omega$ lalu
$\Omega'$ adalah penggeseran sebesar $2\,\vect{\Omega\Omega'}$.
\end{solution}

\begin{solution}{exo:b2:affine:11}
Syarat $\vect{A'B} = \alpha\,\vect{A'C}$ persis mengatakan $1\cdot
\vect{A'B} - \alpha\,\vect{A'C} = 0$, yakni $A' =
\operatorname{bar}(B, 1;\ C, -\alpha)$ (yang berbobot total $1 -
\alpha \neq 0$ karena $B \neq C$); demikian pula $B' =
\operatorname{bar}(C, 1;\ A, -\beta)$ dan $C' =
\operatorname{bar}(A, 1;\ B, -\gamma)$.

\emph{Kriteria kesegarisannya.} Berikanlah tiap titik $P$ baris
barisentrik ternormalkannya $p = (p_A, p_B, p_C)$ dengan $p_A + p_B +
p_C = 1$, terhadap $(A, B, C)$. Bila $\sum_i c_i p_i = 0$
dengan $(c_1, c_2, c_3) \neq 0$ untuk tiga titik $P_1, P_2,
P_3$, maka menjumlahkan entrinya memberi $\sum c_i = 0$, dan
$\sum_i c_i \vect{OP_i} = \sum_j \bigl(\sum_i
c_ip_{ij}\bigr)\vect{OV_j} = 0$: jadi $P_i$ bergantungan secara
\omterm{def:b2:affine:subspace}{afin}, yakni segaris. Sebaliknya sebuah kebergantungan \omterm{def:b2:affine:subspace}{afin}
$(t_i)$ memberi $w = \sum t_ip_i$ yang entrinya berjumlah $0$
dan $\sum_j w_j\vect{OV_j} = 0$; lalu dengan menguraikannya dari $A$: $w_B
\vect{AB} + w_C\vect{AC} = 0$, jadi $w = 0$ berkat kebebasan
\omterm{def:b2:affine:subspace}{afin} $(A, B, C)$: sehingga barisnya bergantungan secara linear.
Karena itu kesegarisannya menyusut menjadi lenyapnya \omterm{def:b2:linalg:det}{determinan} $3 \times 3$,
dan menskalakan barisnya dengan faktor tak nol $1 - \alpha$, $1 -
\beta$, $1 - \gamma$ tak mengubah apa pun:
\[
\det\begin{pmatrix} 0 & 1 & -\alpha\\ -\beta & 0 & 1\\ 1 &
-\gamma & 0\end{pmatrix}
= 1 - \alpha\beta\gamma .
\]
Karena itu $A', B', C'$ segaris bila dan hanya bila $\alpha\beta\gamma = 1$:
itulah teorema Menelaus.
\end{solution}

\begin{solution}{exo:b2:affine:12}
Misalkan $\Delta = \{\lambda \in \R^{n+1} : \lambda_i \geq 0,\
\sum\lambda_i = 1\}$: yang tertutup dan terbatas di $\R^{n+1}$, sehingga
\omterm{def:b2:metric:compact}{kompak}, sedangkan $K^{n+1}$ \omterm{def:b2:metric:compact}{kompak} sebagai hasil kali berhingga. Pemetaan
\[
\Phi \colon \Delta \times K^{n+1} \to \R^n, \qquad
\Phi(\lambda, x_0, \dots, x_n) = \sum_{i=0}^n \lambda_i x_i
\]
bersifat \omterm{def:b2:metric:continuity}{kontinu}, dan \cref{thm:b2:affine:caratheodory} mengatakan
persis bahwa $\operatorname{conv}(K) = \Phi(\Delta \times
K^{n+1})$: yakni peta \omterm{def:b2:metric:continuity}{kontinu} sebuah himpunan \omterm{def:b2:metric:compact}{kompak}
(\cref{thm:b2:metric:compactprops}), sehingga \omterm{def:b2:metric:compact}{kompak}.

Untuk himpunan yang tertutup: ambillah $S = (\R \times \{0\}) \cup \{(0,
1)\}$, yang tertutup di $\R^2$. Sebuah kombinasi \omterm{def:b2:affine:convex}{cembung} yang menaruh bobot
$t$ pada $(0,1)$ dan $1 - t$ pada titik sumbunya berkoordinat kedua
$t$, jadi
\[
\operatorname{conv}(S) = \bigl(\R \times \intco01\bigr) \cup
\{(0,1)\}
\]
(sebab untuk $0 \leq t < 1$ berlaku $(x, t) = t\,(0,1) + (1-t)\,\bigl(\tfrac
x{1-t}, 0\bigr)$). Adapun titik $(1, 1) = \lim_{t \to 1}(1, t)$
bersifat lekat tetapi tak ada di bungkusnya: jadi tak tertutup.
\end{solution}

\begin{solution}{pb:b2:affine:1}
\textbf{1.} Pangkalkan ulang di $A_j$: sebab untuk $i \neq j$, $\vect{A_jA_i} =
\vect{A_0A_i} - \vect{A_0A_j}$. Bila $\sum_{i \neq j}
c_i\vect{A_jA_i} = 0$, maka menguraikannya memberi $\sum_{i \notin \{0,
j\}} c_i\,\vect{A_0A_i} - \bigl(\sum_{i\neq j}
c_i\bigr)\vect{A_0A_j} = 0$; jadi kebebasan
$\vect{A_0A_i}$ memaksa $c_i = 0$ untuk $i \notin \{0, j\}$, lalu
$c_0 = 0$: sehingga kebebasan di $A_j$. Adapun kesetaraannya: dua
keluarga bobot $(\lambda_i)$ dan $(\lambda_i')$ yang berjumlah $1$
dengan \omterm{def:b2:affine:barycenter}{barisentrum} yang sama memberi, dengan $\nu = \lambda -
\lambda'$: $\sum\nu_i = 0$ dan (dengan titik asal $A_0$) $\sum_{i \geq
1}\nu_i\,\vect{A_0A_i} = 0$, jadi $\nu = 0$ di bawah kebebasannya.
Sebaliknya sebuah hubungan tak trivial $\sum_{i\geq1}\mu_i
\vect{A_0A_i} = 0$, yang dilengkapi $\mu_0 = -\sum_{i\geq1}
\mu_i$, memungkinkan kita menambahkan $t(\mu_i)$ pada sembarang keluarga bobot tanpa
menggerakkan \omterm{def:b2:affine:barycenter}{barisentrumnya}: jadi ketaktunggalan.

\textbf{2.} Pasangan $(\vect{AB}, \vect{AC})$ adalah basis $\R^2$:
tulislah $\vect{AM} = \beta\,\vect{AB} + \gamma\,\vect{AC}$
(secara tunggal) lalu tetapkan $\alpha = 1 - \beta - \gamma$; sebab syarat
\omterm{def:b2:affine:barycenter}{barisentrumnya} di titik asal $A$ berbunyi persis $\vect{AM} =
\beta\,\vect{AB} + \gamma\,\vect{AC}$. Adapun ketunggalannya adalah pertanyaan
1.

\textbf{3.} Dari $\alpha\vect{MA} + \beta\vect{MB} +
\gamma\vect{MC} = 0$ dan Chasles, berlaku $\vect{MA} =
-\beta\,\vect{AB} - \gamma\,\vect{AC}$. Dengan $D =
\det(\vect{AB}, \vect{AC})$:
\begin{align*}
\det(\vect{MB}, \vect{MC})
&= \det(\vect{MA} + \vect{AB},\ \vect{MA} + \vect{AC})\\
&= \det(\vect{MA}, \vect{AC}) + \det(\vect{AB}, \vect{MA}) +
D = -\beta D - \gamma D + D = \alpha D,
\end{align*}
dengan memakai kebilinearannya beserta $\det(\vect{MA}, \vect{AC}) = -\beta D$
dan $\det(\vect{AB}, \vect{MA}) = -\gamma D$. Adapun kedua rumus
lainnya menyusul lewat perhitungan yang sama dengan peran yang
dipermutasikan secara \omterm{def:b2:structures:generated}{siklik}.

\textbf{4.} Menurut definisinya, $\operatorname{conv}\{A, B, C\}$ adalah
himpunan \omterm{def:b2:affine:barycenter}{barisentrum} berbobot taknegatif; jadi setelah menormalkan
bobotnya agar berjumlah $1$ lalu memanggil ketunggalannya (pertanyaan 2),
$M \in \operatorname{conv}\{A,B,C\}$ bila dan hanya bila $\alpha, \beta,
\gamma \geq 0$. Tiap koordinatnya merupakan fungsi \omterm{def:b2:affine:subspace}{afin} atas $M$
(lewat pertanyaan 3: yakni \omterm{def:b2:linalg:det}{determinan} $2\times2$ yang satu kolomnya \omterm{def:b2:affine:subspace}{afin}
dalam $M$), jadi tiap syarat tanda \omterm{def:b2:metric:topology}{terbukanya} mendefinisikan setengah
bidang yang \omterm{def:b2:metric:topology}{terbuka}. Adapun pola $(-,-,-)$ bertentangan dengan $\alpha + \beta
+ \gamma = 1$; sedangkan ketujuh pola sisanya terwujud:
skalakanlah sebuah tripel yang menghormati tandanya dengan sekurangnya satu
entri $+$ agar jumlahnya (yang positif) menjadi $1$ --- misalnya
$(-1, 1, 1)$, $(3, -1, -1)$, $(\frac13, \frac13, \frac13)$,
beserta permutasinya.

\textbf{5.} Sebuah \omterm{def:b2:affine:subspace}{pemetaan afin} memelihara \omterm{def:b2:affine:barycenter}{barisentrum}
(\cref{def:b2:affine:subspace}), jadi $u(M) = \alpha u(A) +
\beta u(B) + \gamma u(C)$. Dengan menulis $u(x, y) = ax + by + c$
dan $(a, b) \neq (0,0)$: maka $\{u = c'\}$ merupakan garis, dan setiap
garis $ax + by = c'$ merupakan himpunan aras semacam itu. Bila $u(M), u(N) \geq
c$ dan $t \in \intcc01$, maka $u\bigl(\operatorname{bar}(M,
1-t; N, t)\bigr) = (1-t)u(M) + tu(N) \geq c$: jadi setengah bidangnya
\omterm{def:b2:affine:convex}{cembung}.

\textbf{6.} Syarat $\sum\mu_i = 0$, $3\mu_2 + \mu_4 =
0$, $3\mu_3 + \mu_4 = 0$ memberi (dengan mengambil $\mu_4 = 3$)
kebergantungan $(\mu_1, \mu_2, \mu_3, \mu_4) = (-1, -1, -1, 3)$.
Tandanya terpecah sebagai $\{A_4\} \mid \{A_1, A_2, A_3\}$, lalu
setelah tiap sisinya dinormalkan dengan $3$:
\[
A_4 = \operatorname{bar}\bigl(A_1, \tfrac13;\ A_2, \tfrac13;\
A_3, \tfrac13\bigr) = (1,1),
\]
yakni titik berat segitiganya: jadi titik Radonnya adalah $A_4$
itu sendiri, yang memang terletak di dalam segitiga $A_1A_2A_3$.

\textbf{7.} Pemetaan linear $\Phi \colon \R^{n+2} \to \R
\times \R^n$ dengan $\mu \mapsto (\sum\mu_i,\
\sum\mu_i\vect{OA_i})$ berrank $\leq n + 1$, jadi $\dim\ker
\Phi \geq 1$. Andaikan dua kebergantungan bebas $\mu, \mu'$
ada, maka kombinasi yang sesuai $\nu = \mu'_{n+2}\mu -
\mu_{n+2}\mu'$ (atau $\mu$ itu sendiri bila kedua koefisien terakhirnya
lenyap) akan menjadi kebergantungan \emph{tak nol} dengan $\nu_{n+2}
= 0$; lalu dengan membatasinya ke $A_1, \dots, A_{n+1}$ dan memangkalkan ulang di
$A_1$, ada $\nu_i \neq 0$ dengan $i \geq 2$ (sebab satu bobot tak nol
saja tak dapat berjumlah nol), yang memberi hubungan tak trivial
$\sum_{i\geq2}\nu_i\vect{A_1A_i} = 0$: jadi $n+1$ titiknya akan
bergantungan secara \omterm{def:b2:affine:subspace}{afin}, yang melawan posisi umumnya. Karena itu $\dim\ker
\Phi = 1$. Adapun hujah pembatasan yang sama menunjukkan bahwa kebergantungan
tak nol tak punya koefisien yang lenyap.

\textbf{8.} Misalkan $\mu \neq 0$ sebuah kebergantungan, dengan $P = \{i :
\mu_i > 0\}$ dan $N = \{i : \mu_i < 0\}$: keduanya tak kosong
(sebab $\sum\mu_i = 0$ dan $\mu \neq 0$) sekaligus menyeluruh (sebab tak ada koefisien
yang nol). Konstruksi Radon (\cref{exo:b2:affine:7})
menghasilkan titik bungkus bersamanya persis dari pemartisian ini.
Karena kebergantungannya tunggal hingga sebuah skalar tak nol
(pertanyaan 7), pasangan tak terurut $\{P, N\}$ --- sehingga
pemartisian Radonnya --- bersifat tunggal.

\textbf{9.} Bloknya tak kosong, jadi tipenya $(1,3)$ atau
$(2,2)$. Tipe $(1,3)$ dengan blok $\{j\}$: titik Radonnya terletak di
$\operatorname{conv}\{A_j\} = \{A_j\}$, jadi $A_j \in
\operatorname{conv}$ ketiga lainnya; dan ia tak dapat terletak pada sebuah
rusuk (sebab tiga titiknya akan segaris, yang melawan posisi
umumnya), jadi $A_j$ berada di dalam segitiganya. Tipe
$(2,2)$ dengan blok $\{i,j\} \mid \{k,l\}$: titik Radonnya $z$
terletak di $\intcc{A_i}{A_j} \cap \intcc{A_k}{A_l}$, dan $z$
bukan titik ujungnya (sebab itu akan menyegariskan tiga titik): jadi kedua
ruasnya bersilangan di sebuah titik dalam. Lebih jauh tak ada titik yang terletak
di bungkus yang lain: sebab pemuatan semacam itu $A_l =
\operatorname{bar}(A_i, \lambda_i)_{i \neq l}$ dengan
$\lambda_i \geq 0$ merupakan kebergantungan \omterm{def:b2:affine:subspace}{afin} berpola tanda
$(+,+,+,-)$, yang berkat ketunggalannya (pertanyaan 8) akan membuat
pemartisiannya $(1,3)$. Jadi pada kasus $(2,2)$ keempat titiknya
berada pada posisi \omterm{def:b2:affine:convex}{cembung} dan ruas yang bersilangan itu adalah
diagonalnya.

\textbf{10.} Persamaan $\mu_2 + \mu_3 = 0$, $\mu_3 +
\mu_4 = 0$, $\sum\mu_i = 0$ memberi kebergantungan $(1, -1, 1,
-1)$: jadi pemartisian $\{(0,0), (1,1)\} \mid \{(1,0), (0,1)\}$, dan
\[
\operatorname{bar}\bigl((0,0), \tfrac12;\ (1,1),
\tfrac12\bigr) = \bigl(\tfrac12, \tfrac12\bigr) =
\operatorname{bar}\bigl((1,0), \tfrac12;\ (0,1),
\tfrac12\bigr) :
\]
titik Radonnya adalah pusat bujur sangkarnya, tempat kedua
diagonalnya bersilangan --- yakni tipe $(2,2)$, seperti diramalkan gambarnya.

\textbf{11.} Radon yang diterapkan pada $x_1, \dots, x_4$ memberi blok
$I \mid J$ beserta titik $z \in \operatorname{conv}\{x_i : i
\in I\} \cap \operatorname{conv}\{x_j : j \in J\}$. Tetapkan $k
\in \{1, \dots, 4\}$, katakanlah $k \in I$. Setiap $j \in J$
memenuhi $j \neq k$, jadi $x_j \in C_k$ berkat pemilihan $x_j \in
\bigcap_{l \neq j}C_l$; lalu karena $C_k$ \omterm{def:b2:affine:convex}{cembung}, $z \in
\operatorname{conv}\{x_j : j \in J\} \subseteq C_k$. Karena $k$
sembarang, maka $z \in C_1 \cap C_2 \cap C_3 \cap C_4$.

\textbf{12.} Secara induktif atas $m$. Untuk $m = 3$ hipotesisnya adalah
kesimpulannya; sedangkan $m = 4$ adalah pertanyaan 11. Misalkan $m \geq 4$,
lalu andaikan pernyataannya bagi $m$ himpunan, dan ambillah $C_1, \dots,
C_{m+1}$ dengan sifat perpotongan bertiganya. Tetapkan $C_m' =
C_m \cap C_{m+1}$, yang \omterm{def:b2:affine:convex}{cembung}. Keluarga $C_1, \dots, C_{m-1},
C_m'$ punya $m$ anggota; adapun tripel yang menghindari $C_m'$ berpotongan berkat
hipotesisnya, sedangkan tripel $\{C_i, C_j, C_m'\}$ punya
perpotongan $C_i \cap C_j \cap C_m \cap C_{m+1}$, yang tak kosong
berkat pertanyaan 11 yang diterapkan pada $C_i, C_j, C_m, C_{m+1}$ (sebab setiap tiga
di antaranya bertemu, menurut hipotesisnya). Kini hipotesis induksinya
menghasilkan sebuah titik bersama keluarga barunya, yakni bagi semua $m+1$
himpunannya.

\textbf{13.} (a) Rusuk tertutup $\intcc PQ$, $\intcc QR$,
$\intcc RP$ sebuah segitiga tak merosot: setiap duanya berbagi sebuah
titik sudut, tetapi titik bersama ketiganya akan terletak di $\intcc
PQ \cap \intcc RP = \{P\}$ sekaligus di $\intcc QR$, yang justru mengecualikan
$P$. (b) Setiap tiga di antara himpunan $S_i = \{x_1, \dots, x_4\}
\setminus \{x_i\}$ menghilangkan tiga dari keempat titiknya, sehingga menyisakan
tepat satu titik bersama; sedangkan perpotongan totalnya menghilangkan setiap
titiknya. Adapun $S_i$ bersifat berhingga, tak \omterm{def:b2:affine:convex}{cembung}: jadi kecembungannya
hakiki. (c) Berhingga banyak $H_k = \intco{k}{+\infty} \times
\R$ berpotongan di $\intco{k_{\max}}{+\infty} \times \R \neq
\emptyset$, namun tak ada titik yang $x \geq k$ untuk setiap $k \in \N$:
jadi bagi keluarga tak hingga, kekompakannya hakiki.

\textbf{14.} Andaikan $\bigcap_{i \in I}K_i = \emptyset$ lalu
tetapkan $i_0$. Setiap $x \in K_{i_0}$ meleset dari suatu $K_i$, jadi
$K_{i_0} \subseteq \bigcup_{i \in I}(\R^2 \setminus K_i)$, yakni sebuah
peliput oleh himpunan \omterm{def:b2:metric:topology}{terbuka} (sebab $K_i$ \omterm{def:b2:metric:compact}{kompak}, sehingga tertutup). Menurut
Borel--Lebesgue (\cref{thm:b2:metric:borellebesgue}), berhingga
banyak sudah cukup: $K_{i_0} \cap K_{i_1} \cap \dots \cap K_{i_N} =
\emptyset$. Tetapi setiap tiga anggota keluarga berhingga \omterm{def:b2:affine:convex}{himpunan
cembung} ini berpotongan, jadi pertanyaan 12 membuat perpotongannya
tak kosong: sehingga kontradiksi.

\textbf{15.} Tetapkan $D_p = \overline D(p, r)$ untuk $p \in S$:
yakni himpunan \omterm{def:b2:metric:compact}{kompak} yang \omterm{def:b2:affine:convex}{cembung}. Untuk $p, q, s \in S$, hipotesisnya
memberi cakram tertutup $\overline D(z, r)$ yang memuat $p, q, s$;
lalu $\norm{\vect{zp}}, \norm{\vect{zq}}, \norm{\vect{zs}}
\leq r$, yakni $z \in D_p \cap D_q \cap D_s$. Jadi menurut Helly
(pertanyaan 12; sebab keluarganya berhingga) ada $c \in
\bigcap_{p \in S}D_p$: sehingga setiap $p \in S$ memenuhi
$\norm{\vect{cp}} \leq r$, jadi $S \subseteq \overline D(c,
r)$.

\textbf{16.} Misalkan $c = x_{\lceil n/2\rceil}$. Sebuah setengah
garis tertutup yang memuat $c$ berbentuk $\intoc{-\infty}{t}$ dengan $t \geq
c$ atau $\intco{t}{+\infty}$ dengan $t \leq c$. Yang pertama
memuat $x_1, \dots, x_{\lceil n/2\rceil}$: jadi sekurangnya
$\lceil n/2\rceil \geq n/2$ titik. Yang kedua memuat
$x_{\lceil n/2\rceil}, \dots, x_n$: yakni tepat $n - \lceil
n/2\rceil + 1 = \floor{n/2} + 1 > n/2$ titik.

\textbf{17.} Berlaku $\abs{A \cap B} = \abs A + \abs B - \abs{A \cup
B} \geq \abs A + \abs B - n > \tfrac{4n}3 - n = \tfrac n3$,
lalu
\[
\abs{A \cap B \cap C} \geq \abs{A \cap B} + \abs C - n >
\tfrac n3 + \tfrac{2n}3 - n = 0 .
\]

\textbf{18.} Perhatikan $m > 2n/3$. Untuk $T_1, T_2, T_3 \subseteq S$
yang berkardinalitas $m$, pertanyaan 17 menyediakan titik $x \in T_1
\cap T_2 \cap T_3$; lalu $x \in \operatorname{conv}(T_i)$ untuk
tiap $i$: jadi setiap tiga anggota $\mathcal F$ bertemu. Keluarganya
berhingga (sebab berhingga banyak himpunan bagian $S$) dan terdiri atas
\omterm{def:b2:affine:convex}{himpunan cembung}, jadi Helly (pertanyaan 12) memberi $c \in
\bigcap_{\abs T = m}\operatorname{conv}(T)$.

\textbf{19.} Andaikan suatu setengah bidang tertutup $H \ni c$
memuat kurang dari $n/3$ titik $S$. Maka komplemennya $U$
adalah setengah bidang \omterm{def:b2:metric:topology}{terbuka} yang \omterm{def:b2:affine:convex}{cembung}, dengan $\abs{S \cap U} > 2n/3$,
sehingga $\abs{S \cap U} \geq \floor{2n/3} + 1 = m$; jadi pilihlah $T
\subseteq S \cap U$ dengan $\abs T = m$. Lalu
$\operatorname{conv}(T) \subseteq U$ berkat kecembungan $U$, jadi
$c \in \operatorname{conv}(T) \subseteq U$: yang bertentangan
dengan $c \in H$. Karena itu setiap setengah bidang tertutup yang memuat $c$
memuat sekurangnya $n/3$ titik: jadi $c$ merupakan titik penyeimbang.

\textbf{20.} Ambillah segitiga sama sisi bersisi $L$ dan
$\varepsilon = L/100$. Misalkan $c$ sembarang titik; kita akan memamerkan
setengah bidang tertutup yang memuat $c$ dan paling banyak $k$ titik.

\emph{Kasus 1: $c$ berjarak kurang dari $L/10$ dari sebuah titik sudut, katakanlah $B$.}
Arah dari $c$ ke $A$ dan ke $C$ menyimpang dari arah
$B \to A$ dan $B \to C$ paling banyak sebesar
$\arcsin\bigl(\tfrac{L/10}{9L/10}\bigr) = \arcsin\tfrac19$,
jadi keduanya bersudut $\leq \tfrac\pi3 + 2\arcsin\tfrac19 <
\tfrac{2\pi}3$. \emph{Kasus 2: $c$ berjarak $> L/10$
dari semua titik sudutnya.} Bila $c$ ada di segitiganya, maka ketiga
senjang sudut di antara arah $u_A, u_B, u_C$ dari $c$
ke titik sudutnya berjumlah $2\pi$, jadi suatu senjangnya $\leq 2\pi/3$;
sedangkan bila $c$ di luarnya, ketiga arahnya terletak pada sebuah setengah
bidang arah yang \omterm{def:b2:metric:topology}{terbuka} dan dua di antaranya bersudut $<
\pi/2$. Jadi pada setiap kasusnya dua arah, katakanlah ke $X$ dan $Y$,
bersudut $\leq 2\pi/3$; lalu misalkan $w$ pembagi dua satuannya,
sehingga $\langle u_X, w\rangle, \langle u_Y, w\rangle \geq
\cos\tfrac\pi3 = \tfrac12$. Untuk sembarang titik $b$ pada cakram
di sekitar $X$:
\[
\langle \vect{cb}, w\rangle \geq
\tfrac12\norm{\vect{cX}} - \varepsilon > 0,
\]
sebab $\norm{\vect{cX}} \geq L/10 > 2\varepsilon$ (dan
demikian pula untuk $Y$): jadi setengah bidang \omterm{def:b2:metric:topology}{terbuka} $\{\langle \vect{cx},
w\rangle > 0\}$ menelan kedua gugusnya. Sedangkan komplemen
tertutupnya memuat $c$ dan paling banyak $k$ titik gugus
ketiganya. Jadi tak ada titik bidang yang mengalahkan $n/3$: sehingga bersama
pertanyaan 19, konstanta titik penyeimbangnya persis $1/3$.

\textbf{21.} Urutkanlah sudutnya; maka yang terbesar, sebut $\theta$,
memenuhi $\theta \geq \pi/3$ (sebab ketiganya berjumlah $\pi$). Bila
$\theta \geq \pi/2$, katakanlah di $R$, misalkan $M$ titik tengah
sisi berhadapannya $\intcc PQ$. Rumus garis beratnya (yakni $\vect{RM}
= \tfrac12(\vect{RP} + \vect{RQ})$, uraikan lalu hilangkan
$\langle\vect{RP}, \vect{RQ}\rangle$ lewat hukum kosinusnya)
memberi
\[
\norm{\vect{RM}}^2 = \tfrac12\norm{\vect{RP}}^2 +
\tfrac12\norm{\vect{RQ}}^2 - \tfrac14\norm{\vect{PQ}}^2 \leq
\tfrac12\norm{\vect{PQ}}^2 - \tfrac14\norm{\vect{PQ}}^2 =
\tfrac14\norm{\vect{PQ}}^2,
\]
dengan memakai $\norm{\vect{PQ}}^2 = \norm{\vect{RP}}^2 +
\norm{\vect{RQ}}^2 - 2\langle\vect{RP}, \vect{RQ}\rangle \geq
\norm{\vect{RP}}^2 + \norm{\vect{RQ}}^2$ (sebab hasil kali dalamnya
$\leq 0$). Jadi cakram bergaris tengah $\intcc PQ$, yang berjari-jari
$\leq \tfrac12 < \tfrac1{\sqrt3}$, memuat ketiga titiknya
(sedangkan tripel segaris yang merosot masuk ke $\theta = \pi$). Bila
$\theta < \pi/2$ segitiganya lancip; jadi menurut hukum sinusnya
jari-jari lingkar luarnya $R_c = a/(2\sin\theta)$ dengan $a \leq 1$
sisi yang berhadapan dengan $\theta$, lalu $\theta \in
\intco{\pi/3}{\pi/2}$ memberi $\sin\theta \geq \sqrt3/2$, jadi
$R_c \leq 1/\sqrt3$: sehingga cakram lingkar luarnya menuntaskannya.

\textbf{22.} Untuk $p \in S$ misalkan $K_p = \overline D(p,
1/\sqrt3)$: yakni \omterm{def:b2:metric:compact}{kompak} dan \omterm{def:b2:affine:convex}{cembung}. Sembarang tiga titik $p, q, s$ milik $S$
berjarak berpasangan $\leq 1$, jadi pertanyaan 21 memberi sebuah
cakram berjari-jari $1/\sqrt3$ yang memuatnya: sehingga pusatnya terletak di
$K_p \cap K_q \cap K_s$. Jadi menurut Helly \omterm{def:b2:metric:compact}{kompak} (pertanyaan 14, sebab
keluarga sembarang diizinkan) ada $c \in \bigcap_{p\in
S}K_p$: sehingga setiap $p \in S$ berjarak $1/\sqrt3$ dari $c$, yakni $S
\subseteq \overline D(c, 1/\sqrt3)$. Inilah teorema Jung
pada bidangnya.

\textbf{23.} Dengan $G$ sebagai titik beratnya, $\sum_i\vect{GA_i} = 0$,
jadi
\[
\sum_i\norm{\vect{OA_i}}^2 = \sum_i\norm{\vect{OG} +
\vect{GA_i}}^2 = 3\norm{\vect{OG}}^2 + 2\Bigl\langle
\vect{OG}, \sum_i\vect{GA_i}\Bigr\rangle +
\sum_i\norm{\vect{GA_i}}^2,
\]
lalu suku tengahnya lenyap: itulah kesamaan Leibniz. Untuk
segitiga sama sisi bersisi $1$, berlaku $\norm{\vect{GA_i}} =
1/\sqrt3$ (yakni dua pertiga tingginya $\sqrt3/2$), jadi
$\sum_i\norm{\vect{GA_i}}^2 = 1$. Bila $\overline D(O, r)$
memuat titik sudutnya, maka $3r^2 \geq
\sum_i\norm{\vect{OA_i}}^2 = 3\norm{\vect{OG}}^2 + 1 \geq 1$:
jadi $r \geq 1/\sqrt3$, dengan kesamaannya memaksa $O = G$ dan ketiga
jaraknya sama dengan $r$ --- yakni cakram lingkar luarnya. Jadi konstanta
Jung $1/\sqrt3$ bersifat tajam.

\textbf{24.} \emph{Helly di $\R^n$: bila $C_1, \dots, C_m$
(dengan $m \geq n + 1$) himpunan bagian \omterm{def:b2:affine:convex}{cembung} $\R^n$ dan setiap $n + 1$
di antaranya berpotongan, maka semuanya berpotongan.} Kasus dasarnya $m = n +
2$: pilihlah $x_i \in \bigcap_{j\neq i}C_j$; lalu lema Radon
(\cref{exo:b2:affine:7}) memecah $x_1, \dots, x_{n+2}$ menjadi
blok $I \mid J$ dengan titik bungkus bersama $z$, dan untuk tiap
$k$, blok yang tak memuat $k$ terdiri atas titik
$C_k$, jadi $z \in C_k$ berkat kecembungannya, persis seperti pada pertanyaan
11. Adapun langkah induksinya untuk $m \geq n + 2$: gantilah $C_m,
C_{m+1}$ dengan $C_m \cap C_{m+1}$; sebab sebuah $(n+1)$-tupel keluarga
barunya yang memuat anggota terpotong itu setara dengan $n + 2$
himpunan lamanya, yang ditangani kasus dasarnya, sedangkan tupel
lainnya tercakup hipotesisnya. Jadi rampungkanlah lewat hipotesis
induksinya.

\textbf{25.} Aljabar linear masuk satu kali: sebab $n + 2$ vektor pada
ruang berdimensi $(n+1)$ berisi pasangan (yakni bobot total dan
posisi terbobot) mestilah bergantungan --- dan itulah kebergantungan
\omterm{def:b2:affine:subspace}{afinnya}. Adapun tanda koefisiennya memecah titiknya
menjadi kedua blok Radonnya lalu mengubah satu hubungan linear menjadi
kesamaan dua \omterm{def:b2:affine:barycenter}{barisentrum} taknegatif. Sedangkan kecembungannya dipakai
persis dua kali: pada langkah Helly (sebab bungkus titik $C_k$
tinggal di $C_k$) dan pada penerapannya (sebab setengah bidang dan
cakramnya \omterm{def:b2:affine:convex}{cembung}). Adapun dimensi bidangnya hanya masuk
lewat bilangan $4 = 2 + 2$ titik yang disuapkan ke Radon, yakni
lewat ``$3 = 2 + 1$'' pada hipotesis Helly; sedangkan selebihnya
bebas dimensi, seperti ditegaskan pertanyaan 24. Di $\R^n$
konstantanya menjadi: bilangan Helly $n + 1$; konstanta titik penyeimbang
$\frac1{n+1}$ (yakni setiap himpunan berhingga punya titik yang setiap setengah
ruang tertutup lewatnya memuat pecahan $\geq
\frac1{n+1}$ darinya); dan jari-jari Jung
$\sqrt{\frac{n}{2(n+1)}}$ bagi himpunan berdiameter $1$ --- yang sama
dengan $1/\sqrt3$ ketika $n = 2$.

\end{solution}
